% =====================================================================
%  Bien dich: pdfLaTeX (mac dinh cua Overleaf) -- khong can chinh gi.
%  Neu muon dung font he thong thi doi compiler sang XeLaTeX va thay
%  ba dong inputenc/fontenc/lmodern ben duoi bang:
%      \usepackage{fontspec}
%      \setmainfont{TeX Gyre Termes}
% =====================================================================
\documentclass[11pt,a4paper]{article}

\usepackage[utf8]{inputenc}
\usepackage[T5]{fontenc}
\usepackage{lmodern}

\usepackage[margin=2.6cm]{geometry}
\usepackage{amsmath}
\usepackage{amssymb}
\usepackage{booktabs}
\usepackage{array}
\usepackage{caption}
\usepackage{microtype}
\usepackage{enumitem}
\usepackage{float}
\usepackage{tikz}
\usetikzlibrary{arrows.meta, positioning}
\usepackage[hidelinks]{hyperref}

\renewcommand{\tablename}{Bảng}
\renewcommand{\figurename}{Hình}
\renewcommand{\abstractname}{Tóm tắt}
\renewcommand{\refname}{Tài liệu tham khảo}
\renewcommand{\contentsname}{Nội dung}
\renewcommand{\appendixname}{Phụ lục}

\captionsetup[table]{skip=6pt}
\captionsetup[figure]{skip=8pt}
\newcolumntype{R}{>{\raggedleft\arraybackslash}p{2.7cm}}

\title{\bfseries Hợp nhất hai tầng cho HRM-PET:\\
định tuyến đa nguồn và phân lớp bổ sung với cơ chế cổng theo biên}
\author{Nguyễn Thiện Trường}
\date{Tháng 8, 2026}

\begin{document}

\begin{center}
{\large\bfseries Phần phương pháp}\\[2pt]
{\small Trích từ báo cáo kỹ thuật, gồm phân tích hạn chế và phương pháp đề
xuất. Phần thiết lập thí nghiệm, kết quả và khảo sát loại bỏ thành phần nằm ở
báo cáo đầy đủ.}
\end{center}

\bigskip

\section{Phân tích hai hạn chế}
\label{sec:analysis}

\subsection{Hạn chế thứ nhất: tín hiệu định tuyến chưa tối ưu}

Chúng tôi xây dựng một đầu phân lớp dựa trên phép chiếu ngẫu nhiên (random
projection, viết tắt là đầu RP) trên chính đặc trưng LoRA của HRM-PET. Đầu RP
cho độ chính xác định tuyến cao hơn TII ở cả ba bộ dữ liệu. Quan trọng hơn,
\emph{hợp} của hai bộ định tuyến vượt xa kết quả mà HRM-PET đạt được sau khi đã
tinh chỉnh bằng DRM và CRM.

\begin{table}[htbp]
\centering
\caption{Độ chính xác định tuyến (\%), seed 42. Hàng cuối biểu thị tỉ lệ mẫu mà
ít nhất một trong hai bộ định tuyến cho kết quả đúng, tức là giới hạn trên có
thể đạt được khi kết hợp hai nguồn.}
\label{tab:routing}
\begin{tabular}{lrrr}
\toprule
Nguồn định tuyến & ImageNet-R & CIFAR-100 & CUB-200 \\
\midrule
TII                        & 63.80 & 81.21 & 92.96 \\
HRM-PET sau DRM và CRM     & 77.79 & 89.82 & 93.21 \\
Đầu RP                     & 77.12 & 89.51 & 94.02 \\
\midrule
Hợp của TII và đầu RP      & \textbf{82.99} & \textbf{92.86} & \textbf{96.26} \\
\bottomrule
\end{tabular}
\end{table}

Hai bộ định tuyến cho kết quả sai trên những tập mẫu khác nhau. Trên ImageNet-R,
đầu RP định tuyến đúng $19.2\,\%$ số mẫu mà TII định tuyến sai. Tính bổ trợ này
là điều kiện cần để việc hợp nhất có ý nghĩa: nếu hai nguồn cùng sai trên một
tập mẫu, hợp của chúng sẽ không cao hơn nguồn mạnh hơn.

\subsection{Hạn chế thứ hai: định tuyến tốt hơn không đồng nghĩa phân lớp tốt hơn}

Sau khi cải thiện tầng định tuyến, một hiện tượng khác xuất hiện: mức tăng của
độ chính xác định tuyến chỉ chuyển hóa thành độ chính xác phân lớp với tỉ lệ
$14$--$21\,\%$, và trên CUB-200 thì hoàn toàn không chuyển hóa.

Nguyên nhân nằm ở kiến trúc. Đầu phân lớp dùng chung trải trên \emph{toàn bộ}
các lớp đã quan sát, nên một mẫu bị định tuyến sai vẫn thường được phân lớp
đúng. Trong trường hợp đó, việc sửa lại đường định tuyến không đem lại lợi ích
nào; ngược lại, định tuyến lại một mẫu vốn đã được phân lớp đúng có thể làm kết
quả trở nên sai.

Tuy nhiên, đầu RP không chỉ là một bộ định tuyến. Bản thân nó là một bộ phân lớp
hoàn chỉnh mà tầng định tuyến đã bỏ qua. Bảng \ref{tab:cls} đo mức bổ trợ giữa
hai bộ phân lớp.

\begin{table}[htbp]
\centering
\caption{Tính bổ trợ ở cấp phân lớp (\%), seed 42. Cột cuối biểu thị tỉ lệ mẫu
mà đầu phân lớp sau định tuyến cho kết quả sai còn đầu RP cho kết quả đúng.}
\label{tab:cls}
\begin{tabular}{lrrrr}
\toprule
Bộ dữ liệu & Đầu sau định tuyến & Đầu RP & Hợp & Phần đầu RP bổ trợ \\
\midrule
ImageNet-R & 74.31 & 70.08 & 78.89 & $+4.58$ \\
CIFAR-100  & 89.97 & 88.76 & 92.40 & $+2.43$ \\
CUB-200    & 86.48 & 87.50 & 90.35 & $+3.87$ \\
\bottomrule
\end{tabular}
\end{table}

Ngay trên ImageNet-R, nơi đầu RP có độ chính xác tổng thể thấp hơn $4$ điểm, nó
vẫn phân lớp đúng $4.58\,\%$ số mẫu mà đầu phân lớp chính không xử lý được.

\section{Phương pháp đề xuất}

\subsection{Nguồn tín hiệu thứ hai}

\subsubsection{Ý tưởng}

Ý tưởng nền tảng là mở rộng không gian biểu diễn trước khi phân lớp tuyến tính.
Trong không gian đặc trưng $768$ chiều, các lớp chồng lấn nhau tới mức một siêu
phẳng không tách được chúng. Sau khi chiếu lên không gian $10^4$ chiều và áp một
hàm phi tuyến, các lớp có thêm không gian để tách rời, và một bộ phân lớp tuyến
tính trong không gian mới trở nên đủ mạnh.

Ma trận chiếu hoàn toàn ngẫu nhiên và không được huấn luyện. Đây không phải một
sự đánh đổi mà là điều kiện cần: bất kỳ tham số nào được cập nhật bằng gradient
qua chuỗi tác vụ đều có thể bị ghi đè bởi tác vụ sau, trong khi một ma trận đóng
băng sinh từ seed cố định không bao giờ thay đổi và cũng không cần lưu trữ. Việc
loại bỏ hoàn toàn tham số huấn luyện là tiền đề cho tính chất bất biến thứ tự
trình bày ở Mục \ref{sec:order}.

Nguồn tín hiệu thứ hai gồm một phép chiếu ngẫu nhiên đóng băng, sinh lại được từ
một seed cố định, theo sau là hàm phi tuyến ReLU và một lời giải hồi quy ridge
trên ma trận Gram tích lũy. Không có tham số nào được huấn luyện. Đặc trưng $f$
được trích từ \emph{một} bộ tham số LoRA cố định của HRM-PET, đóng vai trò thích
nghi phiên đầu (first-session adaptation):
\begin{align}
h &= \mathrm{ReLU}(f W), \qquad W \text{ đóng băng, sinh lại từ seed},\\
G &= \sum_n h_n h_n^{\top}, \qquad
C = \sum_n h_n\,\mathrm{onehot}(y_n)^{\top},\\
s &= h^{\top} (G + \lambda I)^{-1} C .
\label{eq:rp}
\end{align}

Hệ thống chỉ lưu hai đại lượng tổng hợp $G$ và $C$; ma trận chiếu $W$ được sinh
lại từ seed. Không có đặc trưng của từng mẫu hay ảnh của các tác vụ trước được
lưu trữ, nên giao thức không lưu mẫu vẫn được bảo toàn.

\begin{table}[H]
\centering
\caption{Ký hiệu trong ba phương trình ở trên. Kích thước ghi theo cấu hình dùng
trong báo cáo: $L = 768$ là số chiều đặc trưng, $M = 10^4$ là số chiều sau khi
chiếu, $K$ là số lớp đã quan sát, $n$ là chỉ số chạy trên các mẫu huấn luyện.}
\label{tab:notation}
\begin{tabular}{@{}llp{7.4cm}@{}}
\toprule
Ký hiệu & Kích thước & Ý nghĩa \\
\midrule
$f$        & $L$          & đặc trưng ảnh, trích từ một bộ tham số LoRA cố định \\
$W$        & $L \times M$ & ma trận chiếu ngẫu nhiên, đóng băng, sinh lại được từ seed \\
$h$        & $M$          & biểu diễn sau khi chiếu và áp ReLU \\
$y_n$      & ---          & nhãn của mẫu thứ $n$ \\
$G$        & $M \times M$ & ma trận Gram, cộng dồn qua mọi mẫu đã quan sát \\
$C$        & $M \times K$ & ma trận tương quan giữa chiều biểu diễn và lớp \\
$\lambda$  & vô hướng     & hệ số điều chuẩn của hồi quy ridge \\
$I$        & $M \times M$ & ma trận đơn vị \\
$s$        & $K$          & điểm số cho từng lớp, đầu ra của nguồn tín hiệu thứ hai \\
\bottomrule
\end{tabular}
\end{table}

Ba phương trình tương ứng với ba bước. Phương trình thứ nhất chiếu đặc trưng lên
không gian $M$ chiều rồi áp hàm phi tuyến. Phương trình thứ hai cộng dồn hai đại
lượng thống kê bậc hai: $G$ ghi lại mức đồng xuất hiện giữa các chiều biểu diễn,
$C$ ghi lại liên hệ giữa chiều biểu diễn và lớp. Cả hai được cập nhật ngay khi
mỗi tác vụ được học và không bao giờ cần đọc lại dữ liệu cũ. Phương trình thứ ba
giải hồi quy ridge, cho ma trận trọng số $W_o = (G + \lambda I)^{-1} C$ và điểm
số $s = h^{\top} W_o$; bước này chỉ thực hiện một lần vào cuối mỗi tác vụ.

\subsubsection{Vai trò bắt buộc của hàm phi tuyến}

Nếu bỏ ReLU khỏi \eqref{eq:rp}, ánh xạ $h = fW$ trở thành tuyến tính. Khi đó, dù
$h$ có $10^4$ chiều, không gian sinh của nó vẫn là một không gian con có số chiều
tối đa bằng hạng của $W$, tức $768$. Hồi quy ridge trên $h$ khi đó tương đương về
mặt toán học với hồi quy ridge trên $f$, và phép chiếu không đóng góp gì ngoài
chi phí bộ nhớ. Chính phép cắt của ReLU mới đưa biểu diễn ra khỏi không gian con
đó và sinh ra các chiều đặc trưng thực sự mới.

\subsubsection{Vai trò của nghịch đảo ma trận Gram}

Nếu chỉ sử dụng $C$ mà bỏ qua $(G + \lambda I)^{-1}$, mô hình quy về một bộ phân
lớp theo nguyên mẫu lớp. Cách làm này không hiệu quả trên đặc trưng ViT vì tồn
tại một số hướng có năng lượng lớn xuất hiện ở mọi lớp, khiến các nguyên mẫu
tương quan mạnh với nhau và phần thông tin phân biệt bị lấn át. Nghịch đảo Gram
thực hiện phép tẩy trắng (whitening): những hướng phổ biến có giá trị lớn trong
$G$ sẽ bị giảm trọng số, để lại các hướng thực sự phân biệt được lớp.

Kết quả thực nghiệm ủng hộ lập luận này: một biến thể tái chấm điểm bằng khoảng
cách Mahalanobis trong không gian $768$ chiều gốc, với cả hiệp phương sai riêng
lẫn hiệp phương sai chung, cho độ chính xác định tuyến $84.3$ so với $93.21$ của
baseline trên CUB-200.

\subsubsection{Tính bất biến thứ tự và hệ quả về quên thảm khốc}
\label{sec:order}

Cả $G$ và $C$ đều là tổng cộng dồn qua các mẫu, và phép cộng có tính giao hoán.
Do đó
\begin{equation}
G = \sum_{t=1}^{T} G_t
\end{equation}
không phụ thuộc vào thứ tự các tác vụ, và lời giải sau tác vụ cuối cùng trùng
khớp với lời giải thu được khi huấn luyện tuyến tính trên toàn bộ dữ liệu cùng
lúc. Nguồn tín hiệu thứ hai vì vậy \emph{không thể} quên theo nghĩa của quên thảm
khốc: đây là một tính chất đại số, không phải một kết quả thực nghiệm.

\subsubsection{Ràng buộc về bộ trích đặc trưng}
\label{sec:fixed-extractor}

Đặc trưng $f$ bắt buộc phải đến từ một bộ trích \emph{cố định trong suốt chuỗi
tác vụ}. Có hai lý do độc lập.

Thứ nhất, $G$ và $C$ được cộng dồn qua mọi tác vụ. Phép cộng chỉ có ý nghĩa khi
các số hạng cùng nằm trong một không gian đặc trưng. Nếu mỗi tác vụ sử dụng một
bộ tham số LoRA khác nhau, các vector $h$ sẽ thuộc những không gian không tương
thích và ma trận Gram trở nên vô nghĩa. Hàm phi tuyến ReLU làm vấn đề trầm trọng
hơn: cùng một ảnh qua hai bộ điều hợp khác nhau sẽ kích hoạt hai tập chiều khác
nhau, tức là một thay đổi rời rạc trong cấu trúc vector chứ không phải một sai
lệch nhỏ và đều.

Thứ hai, ràng buộc nhân quả. Thống kê của tác vụ đầu tiên phải được tích lũy ngay
tại thời điểm học tác vụ đó, vì giao thức không lưu mẫu không cho phép truy cập
lại dữ liệu cũ. Tại thời điểm ấy, các bộ tham số LoRA của những tác vụ sau chưa
tồn tại.

Hai ràng buộc kết hợp lại chỉ để lại hai lựa chọn: mạng xương sống chưa thích
nghi, hoặc bộ tham số LoRA của tác vụ đầu tiên. Chúng tôi chọn phương án thứ hai;
Số liệu so sánh giữa hai phương án nằm ở phần khảo sát các hướng đã loại bỏ của báo cáo đầy đủ.

\subsubsection{Đối chiếu với TII}
\begin{table}[htbp]
\centering
\caption{Đối chiếu hai bộ định tuyến. Hai nguồn khai thác hai loại bằng chứng
khác nhau về cùng một câu hỏi, đây là cơ sở để kỳ vọng chúng cho kết quả sai trên
những tập mẫu khác nhau.}
\label{tab:tii-vs-rp}
\begin{tabular}{@{}p{3.3cm}p{4.9cm}p{5.2cm}@{}}
\toprule
 & TII & Đầu RP \\
\midrule
Bằng chứng sử dụng
  & prompt của từng tác vụ
  & đặc trưng ảnh $768$ chiều \\
Đại lượng đo
  & độ bất định của mô hình khi gắn từng bộ prompt
  & mức tương đồng giữa đặc trưng và nguyên mẫu lớp \\
Cách xác định tham số
  & lan truyền ngược qua nhiều epoch
  & lời giải dạng đóng, một bước \\
Tham số được huấn luyện
  & có
  & không \\
Chịu ảnh hưởng của quên
  & có
  & không, do tính bất biến thứ tự \\
Acc@task, ImageNet-R
  & $63.80$
  & $77.12$ \\
\bottomrule
\end{tabular}
\end{table}

Điểm khác biệt cốt lõi nằm ở hàng thứ hai: TII đánh giá mức độ tự tin của mô hình
khi gắn từng bộ prompt, còn đầu RP đánh giá mức tương đồng giữa đặc trưng của mẫu
và nguyên mẫu của từng lớp. Hai tiêu chí khác nhau về bản chất, nên không có lý
do để chúng cùng thất bại trên một tập mẫu.

\subsection{Tầng thứ nhất: hợp nhất ở cấp định tuyến}
\label{sec:stage1}

\subsubsection{Quy đổi điểm số theo lớp thành điểm số theo tác vụ}

Cả hai nguồn đều xuất ra điểm số theo \emph{lớp}. Để dùng cho định tuyến, điểm số
của một tác vụ được lấy bằng giá trị lớn nhất trong các lớp thuộc tác vụ đó:
\begin{equation}
\ell^{\mathrm{task}}_t = \max_{c \,\in\, \mathcal{C}_t} \ell_c ,
\qquad
s^{\mathrm{task}}_t = \max_{c \,\in\, \mathcal{C}_t} s_c ,
\label{eq:taskscore}
\end{equation}
với $\mathcal{C}_t$ là tập lớp của tác vụ $t$, $\ell$ là logit theo lớp của TII
và $s$ là điểm số theo lớp của đầu RP. Cùng một phép rút gọn được áp dụng cho cả
hai nguồn; hai vector thu được chính là $\ell^{\mathrm{TII}}$ và
$s^{\mathrm{RP}}$ xuất hiện trong \eqref{eq:route}. Phép lấy cực đại phù hợp với câu
hỏi cần trả lời --- \emph{tác vụ này có chứa lớp đúng hay không} --- trong khi
phép lấy tổng hoặc trung bình sẽ ưu tiên những tác vụ có nhiều lớp cùng đạt mức
tương đồng trung bình hơn là một lớp đạt mức tương đồng cao.

\subsubsection{Hợp nhất và chuẩn hóa}

Ký hiệu $z(\cdot)$ là phép chuẩn hóa về trung bình $0$ và phương sai $1$, tính
riêng cho từng mẫu trên tập các tác vụ đã quan sát. Điểm số theo tác vụ của hai
nguồn được kết hợp tuyến tính; tác vụ có điểm cao nhất được chọn, và kết quả
định tuyến này được đưa vào các cơ chế DRM, CRM cùng đầu phân lớp sẵn có của
HRM-PET:
\begin{equation}
\mathrm{fused} = w\, z(\ell^{\mathrm{TII}}) + (1-w)\, z(s^{\mathrm{RP}}),
\qquad w = 0.7 .
\label{eq:route}
\end{equation}

Phép chuẩn hóa ở đây là bắt buộc. Logit của TII có độ phân tán lớn hơn nhiều bậc
so với điểm số hồi quy ridge. Nếu kết hợp trực tiếp, TII sẽ chi phối kết quả ở
mọi giá trị trọng số, và độ chính xác định tuyến giảm xuống \emph{thấp hơn} cả
mức của riêng đầu RP: $68.24$ so với $77.12$. Hiện tượng kết quả hợp nhất thấp
hơn từng nguồn thành phần là dấu hiệu đặc trưng của sai lệch thang đo.

\subsubsection{Lựa chọn trọng số hợp nhất}
\label{sec:choose-w}

Trọng số $w$ đặt trên TII, nên $w = 1.0$ tương ứng với baseline. Giá trị $w = 0.7$
được chọn theo tiêu chí \emph{mọi chỉ số cải thiện trên mọi seed}, chứ không theo
tiêu chí cực đại hóa Acc@1 trung bình.

\begin{table}[htbp]
\centering
\caption{Cơ sở lựa chọn $w$, đo trên ImageNet-R với bốn seed, ở giai đoạn chỉ có
tầng hợp nhất định tuyến.}
\label{tab:choose-w}
\begin{tabular}{@{}p{5.4cm}p{4.6cm}p{2.6cm}@{}}
\toprule
 & $w = 0.6$ & $w = 0.7$ \\
\midrule
Acc@1 trung bình & $+0.48$ & $+0.31$ \\
Cải thiện đồng thời cả sáu chỉ số trên mọi seed
  & không (Forgetting và Backward suy giảm ở seed 43) & có \\
\bottomrule
\end{tabular}
\end{table}

Cấu hình $w = 0.6$ cho mức tăng Acc@1 trung bình cao hơn nhưng không ổn định trên
toàn bộ các seed. Chúng tôi ưu tiên tính nhất quán, vì một phương pháp có mức
tăng trung bình cao hơn nhưng thỉnh thoảng làm suy giảm chỉ số ghi nhớ khó được
xem là cải tiến đáng tin cậy.

\subsection{Tầng thứ hai: hợp nhất ở cấp phân lớp với cơ chế cổng}
\label{sec:stage2}

Phiên bản đầu tiên kết hợp hai bộ phân lớp bằng một trọng số $\beta$ cố định cho
mọi mẫu. Cách làm này gây ra một suy giảm có tính hệ thống: giá trị hàm mất mát
entropy chéo (cross-entropy) trên CIFAR-100 tăng $+0.019 \pm 0.003$, dù độ chính
xác vẫn được cải thiện.

Nguyên nhân là trọng số cố định xử lý như nhau hai loại mẫu có bản chất khác
hẳn: mẫu mà đầu phân lớp chính dự đoán với xác suất $p = 0.97$, và mẫu mà nó
chưa phân định rõ giữa các lớp. Trên CIFAR-100, loại thứ nhất chiếm đa số. Việc
dành $30\,\%$ trọng số quyết định cho nguồn thứ hai trên chính những mẫu đó chỉ
làm phẳng phân bố xác suất tại một dự đoán vốn đã đúng: cross-entropy tăng ngay,
trong khi kết quả $\arg\max$ không thay đổi.

Giải pháp là cho $\beta$ thay đổi theo từng mẫu, tỉ lệ nghịch với mức độ chắc
chắn của đầu phân lớp chính. Đại lượng dùng để đo mức độ chắc chắn là biên
(margin) giữa hai lớp có điểm cao nhất của chính nó. Cơ sở của lựa chọn này:
việc hợp nhất chỉ có thể thay đổi dự đoán bằng cách thay thế lớp xếp thứ nhất,
mà ứng viên thay thế khả dĩ nhất là lớp xếp thứ hai. Gọi $p_{(1)}$ và $p_{(2)}$
là hai xác suất lớn nhất của đầu phân lớp sau định tuyến:
\begin{equation}
\beta_i = \beta \left( 1 - \frac{p_{(1)} - p_{(2)}}{p_{(1)} + p_{(2)}} \right),
\qquad \beta = 0.5 .
\label{eq:gate}
\end{equation}

Biên được chuẩn hóa theo tổng của hai xác suất lớn nhất chứ không theo $1$. Đặt
$r = p_{(2)}/p_{(1)}$, biểu thức trong ngoặc trở thành $1 - (1-r)/(1+r)$, tức
\emph{chỉ phụ thuộc vào tỉ số giữa hai lớp dẫn đầu} và độc lập với phần khối
lượng xác suất mà chúng chiếm. Đây là tính chất mong muốn: mức độ đe dọa của lớp
xếp thứ hai được đo bằng vị thế tương đối của nó, không bằng giá trị tuyệt đối.

Điểm số cuối cùng được kết hợp rồi đưa trở lại thang đo của logit ban đầu:
\begin{align}
m_i &= (1 - \beta_i)\, z(\ell_i) + \beta_i\, z(s_i),\\
\hat{\ell}_i &= m_i \cdot \mathrm{std}(\ell_i) + \mathrm{mean}(\ell_i).
\label{eq:affine}
\end{align}

\begin{table}[htbp]
\centering
\caption{Trọng số hiệu dụng của cơ chế cổng theo \eqref{eq:gate} với
$\beta = 0.5$.}
\label{tab:gate-example}
\begin{tabular}{@{}lrrrr@{}}
\toprule
Trạng thái của đầu phân lớp chính & $p_{(1)}$ & $p_{(2)}$ & Biên & $\beta_i$ \\
\midrule
Dự đoán rất chắc chắn   & 0.97 & 0.01 & 0.980 & 0.010 \\
Dự đoán chắc chắn       & 0.90 & 0.02 & 0.957 & 0.022 \\
Chưa phân định rõ       & 0.35 & 0.30 & 0.077 & 0.462 \\
\bottomrule
\end{tabular}
\end{table}

Bảng \ref{tab:gate-example} cho thấy cơ chế hoạt động đúng như thiết kế: trên các
mẫu đã được dự đoán chắc chắn, nguồn thứ hai chỉ giữ khoảng $2\,\%$ quyền quyết
định; trên các mẫu chưa phân định, tỉ lệ này lên tới $46\,\%$. Cơ chế cổng không
đưa thêm siêu tham số nào và áp dụng một quy tắc duy nhất cho cả ba bộ dữ liệu.

Phép biến đổi \eqref{eq:affine} là biến đổi affine theo từng mẫu với hệ số nhân
dương, do đó đơn điệu tăng: thứ tự xếp hạng giữa các lớp được bảo toàn và mọi
chỉ số độ chính xác không thay đổi. Vai trò duy nhất của nó là đưa thang đo của
logit về mức so sánh được với baseline khi tính cross-entropy. Việc chọn mô-men
của $\ell$ làm mốc còn bảo đảm rằng tại $\beta_i = 0$, phương pháp trả về đúng
logit gốc.

\end{document}
